\documentclass{article}
\usepackage{amsmath} % For \text{}
\usepackage{amssymb} % Not strictly needed for this content, but good practice for math

\begin{document}

\[ \text{(a) } \cos(x) + 2\sin(x)\tan(x) = 1 \quad [0; 2\pi[ \]
\[ \Leftrightarrow \cos(x) + 2\sin(x) \frac{\cos(x)}{\sin(x)} = 1 \]
\[ \Leftrightarrow \cos(x) + 2\cos(x) = 1 \]
\[ \Leftrightarrow 3\cos(x) = 1 \]
\[ \Leftrightarrow \cos(x) = \frac{1}{3} \]
\[ \Leftrightarrow \text{arccos}(\frac{1}{3}) = x \]

\[ \text{(b)(1) } \sin(3x) = \sin(\alpha+40^\circ) \quad [0; 360^\circ[ \]
\[ 3x = \alpha + 40^\circ + k \cdot 2\pi \]
\[ 2\alpha = 40^\circ + k \cdot 2\pi \]
\[ k=0: 2\alpha = 40^\circ + 0 \Rightarrow \alpha = 20^\circ \]
\[ k=1: 2\alpha = 40^\circ + 2\pi \Rightarrow \alpha = \frac{40^\circ + 360^\circ}{2} = 200^\circ = \alpha \]
\[ k=2: 2\alpha = 40^\circ + 4\pi \Rightarrow \alpha = \frac{40^\circ + 720^\circ}{2} = 380^\circ \text{ y hors domaine} \]

\[ \text{(2) } 3\alpha = \pi - (\alpha+40^\circ) + k \cdot 2\pi \]
\[ \Leftrightarrow 3x = \pi - \alpha - 40^\circ + k \cdot 2\pi \]
\[ \Leftrightarrow 4\alpha = \pi - 40^\circ + k \cdot 2\pi \]
\[ k=0: 4\alpha = 180^\circ - 40^\circ + 0 \Leftrightarrow 4\alpha = 140^\circ \Leftrightarrow \alpha = \frac{140^\circ}{4} \]
\[ k=1: 4\alpha = 140^\circ + 2 \cdot 180^\circ \Leftrightarrow \alpha = \frac{140^\circ + 2 \cdot 180^\circ}{4} = \frac{140^\circ + 360^\circ}{4} = \frac{500^\circ}{4} = 125^\circ \alpha \]
\[ k=2: \alpha = \frac{140^\circ + k \cdot 2\pi}{4} \Leftrightarrow \alpha = \frac{140^\circ + 4 \cdot 180^\circ}{4} = \frac{140^\circ + 720^\circ}{4} = \frac{860^\circ}{4} = 200^\circ \]
\[ k=3: \alpha = \frac{140^\circ + 6 \cdot 180^\circ}{4} = \frac{140^\circ + 1060^\circ}{4} = \frac{1200^\circ}{4} = 300^\circ \]
\[ k=4: x = \frac{140^\circ + 8 \cdot 180^\circ}{4} = \text{ hors domaine} \]

\[ S = \{20^\circ, 200^\circ, \frac{140^\circ}{4}, 125^\circ, \frac{860^\circ}{4}, 305^\circ\} \]

\end{document}